\chapter{Overview}
\label{chap:overview}

% archon:covers MR4628606PixtonsFormulaAndAbelJacobiTheoryOnThePicardStack.lean MR4628606PixtonsFormulaAndAbelJacobiTheoryOnThePicardStack/Basic.lean

This chapter is the complete blueprint for the formalization of
Bae--Holmes--Pandharipande--Schmitt--Schwarz,
\textit{Pixton's formula and Abel-Jacobi theory on the Picard stack}
(arXiv:2004.08676, MR4628606).
All declarations are organized from foundational combinatorial definitions
through the main theorems.

\section{Introduction}

\subsection{Double ramification cycles}

\begin{theorem}[Existence of the universal twisted DR cycle]
  \label{thm:dr_existence}
  \lean{MR4628606.TODO.drExistence}
  \uses{def:picard_stack, def:operational_chow}
  % SOURCE: [Bae--Holmes--Pandharipande--Schmitt--Schwarz], Theorem 1 (read from references/MR4628606-picard-stack.tex)
  % SOURCE QUOTE: "Let $g\geq 0$ and $d\in \mathbb{Z}$.
% Let $A=(a_1,\ldots,a_n)$ be a vector of integers satisfying
% $$\sum_{i=1}^n a_i=d\, .$$
% Logarithmic compactification of the Abel-Jacobi map yields a universal twisted double
% ramification
% cycle
% $$\mathsf{DR}^{\mathsf{op}}_{g,A} \in {\mathsf{CH}}^g_{\mathsf{op}}(\Picabs_{g,n,d})$$
% which is compatible with the standard double ramification
% cycle
% $$\mathsf{DR}_{g,A, \omega^k}\in \mathsf{CH}_{2g-3+n}(\overline{\cM}_{g,n})\, $$
% in case $d=k(2g-2)$ for $k\geq 0$."
  \textit{Source: Bae--Holmes--Pandharipande--Schmitt--Schwarz, Theorem 1.}
  Let $g \geq 0$, $d \in \mathbb{Z}$, and $A = (a_1,\ldots,a_n)$ with $\sum a_i = d$.
  Logarithmic compactification of the Abel-Jacobi map yields a universal twisted double
  ramification cycle
  \[
    \DRop_{g,A} \in \CHop^g(\Picabs_{g,n,d})
  \]
  compatible with the standard double ramification cycle
  $\mathsf{DR}_{g,A,\omega^k} \in \Chow_{2g-3+n}(\oM_{g,n})$
  when $d = k(2g-2)$ for $k \geq 0$.
\end{theorem}

\begin{proof}
  \uses{def:picard_stack, def:operational_chow}
  Constructed via logarithmic geometry using the stack of tropical divisors of Marcus--Wise.
  A partial resolution of the classical Abel-Jacobi map over $\Picabs_{g,n,d}$ yields the
  operational class.  Compatibility with canonically twisted DR cycles follows from
  Holmes's infinitesimal criterion (Holmes 2019).
\end{proof}

\begin{definition}[Picard stack]
  \label{def:picard_stack}
  \lean{MR4628606.TODO.PicardStack}
  % SOURCE: [Bae--Holmes--Pandharipande--Schmitt--Schwarz], Section 0.1 (read from references/MR4628606-picard-stack.tex)
  % SOURCE QUOTE: "An object of $\Picabs_{g,n}$ over $\mathcal{S}$ is
% a flat family $$\pi:\mathcal{C} \rightarrow \mathcal{S}$$ of prestable{\footnote{A prestable $n$-pointed
% curve is a connected nodal curve with markings at distinct nonsingular points.
% For the entire paper, we avoid the $(g,n)=(1,0)$ case because of
% non-affine stabilizers.}}
%     $n$-pointed genus $g$ curves together with a line bundle
%     $$\mathcal{L} \rightarrow \mathcal{C}\, .$$
% The Picard stack $\Picabs_{g,n}$ is an algebraic (Artin) stack
% which is locally of finite type"
  % SOURCE QUOTE (smoothness): "Since $\Picabs_{g,n,d}$ is smooth and
% the morphism $j_{\Gamma_\delta}$ is proper, representable, and lci,
% we obtain an operational Chow class..."
% (Bae--Holmes--Pandharipande--Schmitt--Schwarz, Section 1.2, line discussing boundary maps)
  \textit{Source: Bae--Holmes--Pandharipande--Schmitt--Schwarz, Sections 0.1 and 1.2.}
  The \emph{Picard stack} $\Picabs_{g,n,d}$ is a smooth algebraic (Artin) stack locally of
  finite type over a field $\field$ of characteristic zero.  Its objects over a scheme $S$
  consist of flat families $\pi\colon \mathcal{C} \to S$ of prestable $n$-pointed genus-$g$
  curves together with a line bundle $\mathcal{L} \to \mathcal{C}$ of relative degree $d$.
  There is a decomposition
  $\Picabs_{g,n} = \bigsqcup_{d \in \mathbb{Z}} \Picabs_{g,n,d}$.
  (Axiomatized for formalization purposes.)
\end{definition}

\begin{definition}[Operational Chow group]
  \label{def:operational_chow}
  \lean{MR4628606.TODO.OperationalChow}
  \uses{def:picard_stack}
  % SOURCE: [Bae--Holmes--Pandharipande--Schmitt--Schwarz], Section 1.2 (read from references/MR4628606-picard-stack.tex)
  % SOURCE QUOTE: "Let $\frak Y$ be a locally finite type algebraic stack over ${\field}$. Let $p$ be an integer. A bivariant class $c$ in the $p$th operational Chow group $\CHop^p(\frak Y)$ is a collection of  homomorphisms
% \[c(\phi)^{m}\colon \Chow_m(B) \to \Chow_{m-p}(B)\]
% for all maps $\phi \colon B \to \frak Y$ where $B$ is a scheme of finite type over ${\field}$, and for all integers $m$, compatible with proper pushforward, flat pullback, and Gysin homomorphisms for regular embeddings (conditions (C1)-(C3) in \cite[Section 17.1]{Fulton1984Intersection-th})."
  \textit{Source: Bae--Holmes--Pandharipande--Schmitt--Schwarz, Section 1.2.}
  Let $\mathfrak{Y}$ be a locally finite type algebraic stack over $\field$.
  A \emph{bivariant class} $c$ in the $p$th \emph{operational Chow group}
  $\CHop^p(\mathfrak{Y})$ is a collection of homomorphisms
  \[
    c(\phi)^m \colon \Chow_m(B) \to \Chow_{m-p}(B)
  \]
  for all maps $\phi\colon B \to \mathfrak{Y}$ with $B$ a finite-type scheme over $\field$
  and all integers $m$, compatible with proper pushforward, flat pullback, and Gysin
  homomorphisms for regular embeddings.
  (Axiomatized for formalization purposes.)
\end{definition}

\begin{definition}[Boundary stratum and gluing map]
  \label{def:boundary_map_j}
  \lean{MR4628606.TODO.BoundaryStratumMap}
  \uses{def:picard_stack, def:prestable_graph_degree, def:operational_chow}
  % SOURCE: [Bae--Holmes--Pandharipande--Schmitt--Schwarz], Section 1.2 (read from references/MR4628606-picard-stack.tex)
  % SOURCE QUOTE: "For $\Gamma_\delta\in \mathsf{G}_{g,n,k}$, let
% $\Picabs_{\Gamma_\delta}$ be the Picard stack,
% $$\epsilon:\Picabs_{\Gamma_\delta} \rightarrow \mathfrak{M}_{\Gamma}\, ,$$
% parameterizing curves with degenerations forced by $\Gamma$ and with  line bundles which
% have
% degree $\delta(v)$ restriction to the components corresponding to the vertex $v\in \V$.
% We have a
% canonical map
% $$j_{\Gamma_\delta} : \Picabs_{\Gamma_\delta} \to \Picabs_{g,n,d}\, .$$
% Since $\Picabs_{g,n,d}$ is smooth and
% the morphism $j_{\Gamma_\delta}$ is proper, representable, and lci,
% we obtain an operational Chow class,
% $$j_{\Gamma_\delta *}[\Picabs_{\Gamma_\delta}] \in \mathsf{CH}_{\mathsf{op}}^{|\E(\Gamma)|}(\Picabs_{g,n,d})\, ."
  \textit{Source: Bae--Holmes--Pandharipande--Schmitt--Schwarz, Section 1.2.}
  For each $\Gamma_\delta \in \GG_{g,n,d}$, the \emph{boundary stratum}
  $\Picabs_{\Gamma_\delta}$ is the Picard stack parameterizing prestable curves
  whose degenerations are forced by $\Gamma$ and whose line bundle has multidegree $\delta$:
  i.e., its restriction to the component corresponding to $v \in \V$ has degree $\delta(v)$.
  The \emph{boundary stratum map} is the canonical morphism
  \[
    j_{\Gamma_\delta}\colon \Picabs_{\Gamma_\delta} \longrightarrow \Picabs_{g,n,d}.
  \]
  The morphism $j_{\Gamma_\delta}$ is proper, representable, and locally complete intersection.
  For a decorated prestable graph $[\Gamma_\delta, \gamma] \in \DGG_{g,n,d}$, the
  \emph{tautological class}
  \[
    j_{\Gamma_\delta*}[\gamma] \in \CHop^*(\Picabs_{g,n,d})
  \]
  is the operational Chow class obtained via pushforward along $j_{\Gamma_\delta}$ and
  the action of the decoration data $\gamma$.
  (Axiomatized for formalization purposes.)
\end{definition}

\subsection{Pixton's formula}

\subsubsection{Prestable graphs}

\begin{definition}[Prestable graph]
  \label{def:prestable_graph}
  \lean{MR4628606.TODO.PrestableGraph}
  % SOURCE: [Bae--Holmes--Pandharipande--Schmitt--Schwarz], Section 0.2.1 (read from references/MR4628606-picard-stack.tex)
  % SOURCE QUOTE: "We define the set $\G_{g,n}$ of {\em prestable graphs} as
% follows. A prestable graph $\Gamma \in \G_{g,n}$ consists of the data
% $$\Gamma=(\V\, ,\ \H\, ,\ \L\, , \ \mathrm{g}:\V \rarr \Z_{\geq 0}\, ,
% \ v:\H\rarr \V\, ,
% \ \iota : \H\rarr \H)$$
% satisfying the properties:
% \begin{enumerate}
% \item[(i)] $\V$ is a vertex set with a genus function $\g:\V\to \Z_{\geq 0}$,
% \item[(ii)] $\H$ is a half-edge set equipped with a
% vertex assignment $v:\H \to \V$ and an involution $\iota$,
% \item[(iii)] $\E$, the edge set, is defined by the
% 2-cycles of $\iota$ in $\H$ (self-edges at vertices
% are permitted),
% \item[(iv)] $\L$, the set of legs, is defined by the fixed points of $\iota$ and is
% placed in bijective correspondence with a set of $n$ markings,
% \item[(v)] the pair $(\V,\E)$ defines a {\em connected} graph
% satisfying the genus condition
% $$\nice \sum_{v \in \V} \g(v) + h^1(\Gamma) = g\, ."
  \textit{Source: Bae--Holmes--Pandharipande--Schmitt--Schwarz, Section 0.2.1.}
  A \emph{prestable graph} $\Gamma \in \GG_{g,n}$ consists of data
  \[
    \Gamma = (\V,\; \HH,\; \LL,\; \mathrm{g}\colon \V \to \mathbb{Z}_{\geq 0},\;
              v\colon \HH \to \V,\; \iota\colon \HH \to \HH)
  \]
  satisfying:
  \begin{enumerate}
  \item[(i)] $\V$ is a finite vertex set with genus function $\mathrm{g}\colon \V \to \mathbb{Z}_{\geq 0}$;
  \item[(ii)] $\HH$ is a finite half-edge set with vertex assignment $v\colon \HH \to \V$
    and an involution $\iota\colon \HH \to \HH$;
  \item[(iii)] the edge set $\E$ consists of the $2$-cycles of $\iota$ (self-edges permitted);
  \item[(iv)] the leg set $\LL$ consists of the fixed points of $\iota$, in bijection with
    $\{1,\ldots,n\}$;
  \item[(v)] the graph $(\V,\E)$ is connected and satisfies the genus condition
    $\sum_{v \in \V} \mathrm{g}(v) + h^1(\Gamma) = g$.
  \end{enumerate}
\end{definition}

\begin{definition}[Isomorphism of prestable graphs]
  \label{def:prestable_graph_isom}
  \lean{MR4628606.TODO.PrestableGraphIsom}
  \uses{def:prestable_graph}
  % SOURCE: [Bae--Holmes--Pandharipande--Schmitt--Schwarz], Section 0.2.1 (read from references/MR4628606-picard-stack.tex)
  % SOURCE QUOTE: "An isomorphism between $\Gamma$, $\Gamma' \in \G_{g,n}$
% consists of bijections $\V \to \V'$ and $\H \to \H'$ respecting the
% structures $\L$, $\mathrm{g}$, $v$, and $\iota$.
% Let $\Aut(\Gamma)$ denote the automorphism group of $\Gamma$."
  \textit{Source: Bae--Holmes--Pandharipande--Schmitt--Schwarz, Section 0.2.1.}
  An \emph{isomorphism} between $\Gamma, \Gamma' \in \GG_{g,n}$ consists of bijections
  $\V \to \V'$ and $\HH \to \HH'$ respecting the structures $\LL$, $\mathrm{g}$, $v$,
  and $\iota$. The \emph{automorphism group} $\Aut(\Gamma)$ is the group of automorphisms
  of $\Gamma$.
\end{definition}

\begin{lemma}[Finiteness of isomorphism classes]
  \label{lem:prestable_graph_finite_isom_classes}
  \lean{MR4628606.TODO.prestableGraphFiniteIsomClasses}
  \uses{def:prestable_graph, def:prestable_graph_isom}
  % SOURCE: [Bae--Holmes--Pandharipande--Schmitt--Schwarz], Section 0.2.1 (read from references/MR4628606-picard-stack.tex)
  % SOURCE QUOTE: "While the set of isomorphism classes of prestable graphs is infinite, the set of isomorphism classes of prestable graphs with prescribed bounds on the number of edges is finite."
  \textit{Source: Bae--Holmes--Pandharipande--Schmitt--Schwarz, Section 0.2.1.}
  For any $N \geq 0$, there are only finitely many isomorphism classes of prestable graphs
  $\Gamma \in \GG_{g,n}$ with $|\E(\Gamma)| \leq N$.
\end{lemma}

\begin{proof}
  \uses{def:prestable_graph, def:prestable_graph_isom}
  The genus condition $\sum_{v} \mathrm{g}(v) + h^1(\Gamma) = g$ together with the
  edge bound $|\E(\Gamma)| \leq N$ constrains the number of vertices, genus labels, and
  total number of half-edges to finitely many possibilities.  The bijection of $\LL$ with
  $\{1,\ldots,n\}$ adds no further parameters.  Hence there are only finitely many
  isomorphism classes.
\end{proof}

\subsubsection{Prestable graphs with degrees}

\begin{definition}[Prestable graph of degree $d$]
  \label{def:prestable_graph_degree}
  \lean{MR4628606.TODO.PrestableGraphDegree}
  \uses{def:prestable_graph}
  % SOURCE: [Bae--Holmes--Pandharipande--Schmitt--Schwarz], Section 0.2.2 (read from references/MR4628606-picard-stack.tex)
  % SOURCE QUOTE: "We define the set $\G_{g,n,d}$ of {\em prestable graphs of degree $d$} as
% follows: $$\Gamma_\delta=(\Gamma,\delta) \in \G_{g,n,d}$$ consists of the data
% \begin{enumerate}
%     \item[$\bullet$] a prestable graph $\Gamma\in \mathsf{G}_{g,n}$,
%     \item[$\bullet$] a function $\delta: \V \rightarrow \mathbb{Z}$
%     satisfying the degree condition
% $$\sum_{v\in V} \delta(v) = d\, .$$
% \end{enumerate}
% The function $\delta$ is  often called the {\em multidegree}."
  \textit{Source: Bae--Holmes--Pandharipande--Schmitt--Schwarz, Section 0.2.2.}
  A \emph{prestable graph of degree $d$} is a pair $\Gamma_\delta = (\Gamma, \delta)$
  where $\Gamma \in \GG_{g,n}$ is a prestable graph and $\delta\colon \V \to \mathbb{Z}$
  is a function (the \emph{multidegree}) satisfying $\sum_{v \in \V} \delta(v) = d$.
  The set of isomorphism classes is denoted $\GG_{g,n,d}$.
\end{definition}

\begin{definition}[Automorphisms of a prestable graph of degree $d$]
  \label{def:aut_prestable_graph_degree}
  \lean{MR4628606.TODO.AutPrestableGraphDegree}
  \uses{def:prestable_graph_degree, def:prestable_graph_isom}
  % SOURCE: [Bae--Holmes--Pandharipande--Schmitt--Schwarz], Section 0.2.2 (read from references/MR4628606-picard-stack.tex)
  % SOURCE QUOTE: "An automorphism of $\Gamma_\delta \in \G_{g,n,d}$
% consists of an automorphism of $\Gamma$
% leaving $\delta$ invariant.
% Let $\Aut(\Gamma_\delta)$ denote the automorphism group of $\Gamma_\delta$."
  \textit{Source: Bae--Holmes--Pandharipande--Schmitt--Schwarz, Section 0.2.2.}
  An \emph{automorphism} of $\Gamma_\delta \in \GG_{g,n,d}$ is an automorphism of the
  underlying prestable graph $\Gamma$ leaving the multidegree $\delta$ invariant.
  The \emph{automorphism group} $\Aut(\Gamma_\delta)$ is finite.
\end{definition}

\subsubsection{Tautological $\psi$, $\xi$, and $\eta$ classes}

\begin{definition}[Tautological classes $\psi$, $\xi$, $\eta$]
  \label{def:tautological_classes}
  \lean{MR4628606.TODO.TautologicalClasses}
  \uses{def:picard_stack, def:operational_chow}
  % SOURCE: [Bae--Holmes--Pandharipande--Schmitt--Schwarz], Definition in Section 0.2.3 (read from references/MR4628606-picard-stack.tex)
  % SOURCE QUOTE: "The following operational classes on $\Picabs_{g,n}$ are obtained from the
% universal structures:
% \begin{itemize}
%     \item $\psi_i = c_1(p_i^* \omega_\pi)\in \mathsf{CH}^1_{\mathsf{op}}(\Picabs_{g,n})\, $,
%     \item $\xi_i = c_1(p_i^* \mathfrak{L})\in \mathsf{CH}^1_{\mathsf{op}}(\Picabs_{g,n})
%     \, $,
%     \item $\eta_{a,b} = \pi_*\left(c_1(\omega_\log)^a \xi^b\right)
%     \in \mathsf{CH}^{a+b-1}_{\mathsf{op}}(\Picabs_{g,n})
%     \, $.
% \end{itemize}"
  \textit{Source: Bae--Holmes--Pandharipande--Schmitt--Schwarz, Definition in Section 0.2.3.}
  The following operational classes on $\Picabs_{g,n}$ are obtained from the universal
  curve $\pi\colon \mathfrak{C}_{g,n} \to \Picabs_{g,n}$, the universal line bundle
  $\mathfrak{L} \to \mathfrak{C}_{g,n}$, and the $i$th section $p_i$:
  \begin{itemize}
  \item $\psi_i = c_1(p_i^* \omega_\pi) \in \CHop^1(\Picabs_{g,n})$,
  \item $\xi_i = c_1(p_i^* \mathfrak{L}) \in \CHop^1(\Picabs_{g,n})$,
  \item $\eta_{a,b} = \pi_*(c_1(\omega_{\log})^a \xi^b) \in \CHop^{a+b-1}(\Picabs_{g,n})$,
  \end{itemize}
  where $\omega_{\log} = \omega_\pi(\sum_{i=1}^n \mathfrak{S}_i)$ is the relative
  log-canonical bundle and $\xi = c_1(\mathfrak{L})$.
  We set $\eta = \eta_{0,2} = \pi_*(\xi^2)$.
\end{definition}

\begin{definition}[Decorated prestable graph]
  \label{def:decorated_prestable_graph}
  \lean{MR4628606.TODO.DecoratedPrestableGraph}
  \uses{def:prestable_graph_degree}
  % SOURCE: [Bae--Holmes--Pandharipande--Schmitt--Schwarz], Section 0.2.3 (read from references/MR4628606-picard-stack.tex)
  % SOURCE QUOTE: "A {\em decorated prestable graph} $[\Gamma_\delta,\gamma]$ {\em of degree $d$}
% is a prestable graph $\Gamma_\delta \in \G_{g,n,d}$ of degree $d$ together with the following
% decoration data $\gamma$:
% \begin{itemize}
% \item each leg $i\in \L$ is decorated with a monomial $\psi_i^a \xi_i^b$,
% \item each half-edge $h\in \H\setminus \L$ is
% decorated with a monomial $\psi^a_h$,
% \item each edge $e\in \E$ is decorated with a monomial $\xi^a_e$,
% \item each vertex in $\V$ is decorated with a monomial in the variables
% $\{\eta_{a,b}\}_{a+b \geq 2}$.
% \end{itemize}
% In all four cases, the monomial may be be trivial."
  \textit{Source: Bae--Holmes--Pandharipande--Schmitt--Schwarz, Section 0.2.3.}
  A \emph{decorated prestable graph of degree $d$} is a pair $[\Gamma_\delta, \gamma]$
  where $\Gamma_\delta \in \GG_{g,n,d}$ and $\gamma$ is \emph{decoration data} consisting of:
  \begin{itemize}
  \item a monomial $\psi_i^a \xi_i^b$ for each leg $i \in \LL$;
  \item a monomial $\psi_h^a$ for each half-edge $h \in \HH \setminus \LL$;
  \item a monomial $\xi_e^a$ for each edge $e \in \E$;
  \item a monomial in $\{\eta_{a,b}\}_{a+b \geq 2}$ for each vertex $v \in \V$.
  \end{itemize}
  Any monomial may be trivial.  The set of decorated prestable graphs of degree $d$
  is denoted $\DGG_{g,n,d}$.
\end{definition}

\begin{definition}[Tautological ring]
  \label{def:tautological_ring}
  \lean{MR4628606.TODO.TautologicalRing}
  \uses{def:tautological_classes, def:decorated_prestable_graph,
        def:operational_chow, def:picard_stack, def:boundary_map_j}
  % SOURCE: [Bae--Holmes--Pandharipande--Schmitt--Schwarz], Section 0.2.3 (read from references/MR4628606-picard-stack.tex)
  % SOURCE QUOTE: "The {\em tautological  classes} in
%  $\mathsf{CH}^*_{\mathsf{op}}(\Picabs_{g,n,d})$ consist of the
%  $\mathbb{Q}$-linear span
%  of the operational classes associated to
%  all $[\Gamma_\delta,\gamma]\in \D\G_{g,n,d}$."
  \textit{Source: Bae--Holmes--Pandharipande--Schmitt--Schwarz, Section 0.2.3.}
  The \emph{tautological classes} in $\CHop^*(\Picabs_{g,n,d})$ are the
  $\mathbb{Q}$-linear span of the operational classes $j_{\Gamma_\delta*}[\gamma]$
  associated to all decorated prestable graphs $[\Gamma_\delta, \gamma] \in \DGG_{g,n,d}$.
\end{definition}

\subsubsection{Weightings mod $r$}

\begin{definition}[Weighting mod $r$]
  \label{def:weighting_mod_r}
  \lean{MR4628606.TODO.WeightingModR}
  \uses{def:prestable_graph_degree}
  % SOURCE: [Bae--Holmes--Pandharipande--Schmitt--Schwarz], Section 0.2.4 (read from references/MR4628606-picard-stack.tex)
  % SOURCE QUOTE: "A {\em weighting mod~$r$} of $\Gamma_\delta$ is a function on the set of half-edges,
% $$ w:\H(\Gamma_\delta) \rightarrow \{0,1, \ldots, r-1\}\, ,$$
% which satisfies the following three properties:
% \begin{enumerate}
% \item[(i)] $\forall i\in \L(\Gamma_\delta)$, corresponding to
%  the marking $i\in \{1,\ldots, n\}$,
% $$w(i)=a_i  \mod r\, ,$$
% \item[(ii)] $\forall e \in \E(\Gamma_\delta)$, corresponding to two half-edges
% $h,h' \in \H(\Gamma_\delta)$,
% $$w(h)+w(h')=0 \mod r\, ,$$
% \item[(iii)] $\forall v\in \V(\Gamma_\delta)$,
% $$\sum_{v(h)= v} w(h)= \delta(v) \mod r\, ,$$
% where the sum is taken over {\em all} $\mathsf{n}(v)$ half-edges incident
% to $v$.\end{enumerate}"
  \textit{Source: Bae--Holmes--Pandharipande--Schmitt--Schwarz, Section 0.2.4.}
  Let $\Gamma_\delta \in \GG_{g,n,d}$, $A = (a_1,\ldots,a_n) \in \mathbb{Z}^n$ with
  $\sum a_i = d$, and $r$ a positive integer.
  A \emph{weighting mod $r$} of $\Gamma_\delta$ is a function
  $w\colon \HH(\Gamma_\delta) \to \{0,1,\ldots,r-1\}$ satisfying:
  \begin{enumerate}
  \item[(i)] For each leg $i \in \LL(\Gamma_\delta)$ corresponding to marking $i \in \{1,\ldots,n\}$:
    $w(i) \equiv a_i \pmod{r}$.
  \item[(ii)] For each edge $e \in \E(\Gamma_\delta)$ with half-edges $h, h'$:
    $w(h) + w(h') \equiv 0 \pmod{r}$.
  \item[(iii)] For each vertex $v \in \V(\Gamma_\delta)$:
    $\sum_{v(h) = v} w(h) \equiv \delta(v) \pmod{r}$,
    where the sum is over all $\mathrm{n}(v)$ half-edges incident to $v$.
  \end{enumerate}
  The set of all weightings mod $r$ of $\Gamma_\delta$ is denoted
  $\mathsf{W}_{\Gamma_\delta,r}$.
\end{definition}

\begin{lemma}[Cardinality of weightings mod $r$]
  \label{lem:weightings_cardinality}
  \lean{MR4628606.TODO.weightingsCardinality}
  \uses{def:weighting_mod_r, def:prestable_graph_degree}
  % SOURCE: [Bae--Holmes--Pandharipande--Schmitt--Schwarz], Section 0.2.4 (read from references/MR4628606-picard-stack.tex)
  % SOURCE QUOTE: "We denote by $\mathsf{W}_{\Gamma_\delta,r}$ the finite set of all possible weightings
% mod $r$
% of
% $\Gamma_\delta$. The set $\mathsf{W}_{\Gamma_\delta,r}$ has cardinality $r^{h^1(\Gamma_\delta)}$."
  \textit{Source: Bae--Holmes--Pandharipande--Schmitt--Schwarz, Section 0.2.4.}
  The set $\mathsf{W}_{\Gamma_\delta,r}$ has cardinality $r^{h^1(\Gamma_\delta)}$.
\end{lemma}

\begin{proof}
  \uses{def:weighting_mod_r, def:prestable_graph_degree}
  Condition (i) fixes the values $w(i)$ on legs up to the constraint $a_i \bmod r$.
  Condition (ii) pairs the two half-edges at each edge: once one is chosen freely in
  $\{0,\ldots,r-1\}$, the other is determined.  Condition (iii) at each vertex is then
  redundant except for exactly $h^0(\Gamma_\delta) - 1 = 0$ independent constraints
  (since the graph is connected).  Therefore the free parameters are one per independent
  cycle, giving $r^{h^1(\Gamma_\delta)}$ choices in total.
\end{proof}

\subsubsection{Calculation of the twisted double ramification cycle}

\begin{definition}[Pixton's formula mod $r$]
  \label{def:pixton_formula_r}
  \lean{MR4628606.TODO.PixtonFormulaR}
  \uses{def:prestable_graph_degree, def:weighting_mod_r, def:tautological_classes,
        def:aut_prestable_graph_degree, def:operational_chow, def:picard_stack,
        def:boundary_map_j, def:decorated_prestable_graph}
  % SOURCE: [Bae--Holmes--Pandharipande--Schmitt--Schwarz], Section 0.2.4 (read from references/MR4628606-picard-stack.tex)
  % SOURCE QUOTE: "We denote by
% $\P_{g,A,d}^{c,r}\in \mathsf{CH}^c_{\mathsf{op}}(\Picabs_{g,n,d})$
% the codimension
%  $c$ component of the tautological operational class
% \begin{multline*}
% \hspace{-10pt}\sum_{
% \substack{\Gamma_\delta\in \G_{g,n,d} \\
% w\in \mathsf{W}_{\Gamma_\delta,r}}
% }
% \frac{r^{-h^1(\Gamma_\delta)}}{|\Aut(\Gamma_\delta)| }
% \;
% j_{\Gamma_\delta*}\Bigg[
% \prod_{i=1}^n \exp\left(\frac12 a_i^2 \psi_i + a_i \xi_i \right)
% \prod_{v \in \V(\Gamma_\delta)} \exp\left(-\frac12 \eta(v) \right)
% \\ \hspace{+10pt}
% \prod_{e=(h,h')\in \E(\Gamma_\delta)}
% \frac{1-\exp\left(-\frac{w(h)w(h')}2(\psi_h+\psi_{h'})\right)}{\psi_h + \psi_{h'}} \Bigg]\, ."
  \textit{Source: Bae--Holmes--Pandharipande--Schmitt--Schwarz, Section 0.2.4.}
  For $g \geq 0$, $A = (a_1,\ldots,a_n) \in \mathbb{Z}^n$ with $\sum a_i = d$, and $r$ a
  positive integer, define $\Pix^{c,r}_{g,A,d} \in \CHop^c(\Picabs_{g,n,d})$ as the
  codimension-$c$ component of
  \[
    \sum_{\substack{\Gamma_\delta \in \GG_{g,n,d} \\ w \in \mathsf{W}_{\Gamma_\delta,r}}}
    \frac{r^{-h^1(\Gamma_\delta)}}{|\Aut(\Gamma_\delta)|}
    \; j_{\Gamma_\delta*}\!\left[
      \prod_{i=1}^n e^{\frac{1}{2}a_i^2 \psi_i + a_i \xi_i}
      \cdot \prod_{v \in \V(\Gamma_\delta)} e^{-\frac{1}{2}\eta(v)}
      \cdot \prod_{e=(h,h') \in \E(\Gamma_\delta)}
        \frac{1 - e^{-\frac{w(h)w(h')}{2}(\psi_h + \psi_{h'})}}{\psi_h + \psi_{h'}}
    \right].
  \]
  The sum is over isomorphism classes; only finitely many terms yield a nonzero operation
  on any given family over a finite-type base.
\end{definition}

\begin{proposition}[Polynomiality of Pixton's formula]
  \label{prop:pixton_polynomiality}
  \lean{MR4628606.TODO.pixtonPolynomiality}
  \uses{def:pixton_formula_r, def:decorated_prestable_graph, def:weighting_mod_r}
  % SOURCE: [Bae--Holmes--Pandharipande--Schmitt--Schwarz], Proposition in Section 0.2.4 (read from references/MR4628606-picard-stack.tex)
  % SOURCE QUOTE: "For fixed $g$, $A$, $d$, and $c$ and a decorated graph $[\Gamma_\delta, \gamma]$ of
% degree $d$, the coefficient of $j_{\Gamma_\delta*} [\gamma]$ in the
% tautological class
% $$\P_{g,A,d}^{c,r} \in \mathsf{CH}^c_{\mathsf{op}}(\Picabs_{g,n,d})$$
% is polynomial in $r$ (for all sufficiently large $r$)."
  \textit{Source: Bae--Holmes--Pandharipande--Schmitt--Schwarz, Proposition in Section 0.2.4.}
  For fixed $g$, $A$, $d$, $c$, and a decorated graph $[\Gamma_\delta, \gamma] \in \DGG_{g,n,d}$,
  the coefficient of $j_{\Gamma_\delta*}[\gamma]$ in
  $\Pix^{c,r}_{g,A,d} \in \CHop^c(\Picabs_{g,n,d})$
  is polynomial in $r$ for all sufficiently large $r$.
\end{proposition}

% SOURCE QUOTE PROOF: "The following fundamental polynomiality property of $\P_{g,A,d}^{c,r}$
%  is
% parallel to Pixton's polynomiality in \cite[Appendix]{Janda2016Double-ramifica} and
% is a consequence of \cite[Proposition $3''$]{Janda2016Double-ramifica}."
\begin{proof}
  \uses{def:pixton_formula_r, def:weighting_mod_r, lem:weightings_cardinality}
  The set $\mathsf{W}_{\Gamma_\delta,r}$ has cardinality $r^{h^1(\Gamma_\delta)}$ by
  \cref{lem:weightings_cardinality}.  Expanding the exponentials and edge factors in
  the definition of $\Pix^{c,r}_{g,A,d}$ yields, for each decorated graph, a finite sum of
  products of weight values.  Each such product is polynomial in $r$ by the polynomiality
  result of Janda--Pandharipande--Pixton--Zvonkine (Proposition $3''$ of
  \textit{Double ramification cycles on the moduli spaces of curves}).
\end{proof}

\begin{definition}[Pixton's formula]
  \label{def:pixton_formula}
  \lean{MR4628606.TODO.PixtonFormula}
  \uses{prop:pixton_polynomiality}
  % SOURCE: [Bae--Holmes--Pandharipande--Schmitt--Schwarz], Section 0.2.4 (read from references/MR4628606-picard-stack.tex)
  % SOURCE QUOTE: "We denote by $\P_{g,A,d}^c$ the value at $r=0$
% of the polynomial associated to $\P_{g,A,d}^{c,r}$ by Proposition~\ref{pply}. In other words, $\P_{g,A,d}^c$ is the {\em constant} term of the associated polynomial in $r$."
  \textit{Source: Bae--Holmes--Pandharipande--Schmitt--Schwarz, Section 0.2.4.}
  The class $\Pix^c_{g,A,d} \in \CHop^c(\Picabs_{g,n,d})$ is the value at $r = 0$ of
  the polynomial in $r$ given by \cref{prop:pixton_polynomiality}; equivalently,
  $\Pix^c_{g,A,d}$ is the constant term of that polynomial.
\end{definition}

\begin{theorem}[Main theorem: Pixton's formula computes the DR cycle]
  \label{thm:main}
  \lean{MR4628606.TODO.mainTheorem}
  \uses{thm:dr_existence, def:pixton_formula, prop:pixton_polynomiality,
        def:tautological_classes, def:tautological_ring}
  % SOURCE: [Bae--Holmes--Pandharipande--Schmitt--Schwarz], Theorem 2 (read from references/MR4628606-picard-stack.tex)
  % SOURCE QUOTE: "Let $g\geq 0$ and $d\in \mathbb{Z}$.
% Let $A=(a_1,\ldots,a_n)$ be a vector of integers with
% $\sum_{i=1}^n a_i=d\,$.
% The universal twisted double ramification cycle is calculated by Pixton's formula:
% $$\mathsf{DR}^{\mathsf{op}}_{g,A} =\P_{g,A,d}^g\
% \in {\mathsf{CH}}^g_{\mathsf{op}}(\Picabs_{g,n,d})\, .$$"
  \textit{Source: Bae--Holmes--Pandharipande--Schmitt--Schwarz, Theorem 2.}
  Let $g \geq 0$, $d \in \mathbb{Z}$, and $A = (a_1,\ldots,a_n)$ with $\sum a_i = d$.
  The universal twisted double ramification cycle equals Pixton's formula:
  \[
    \DRop_{g,A} = \Pix^g_{g,A,d} \in \CHop^g(\Picabs_{g,n,d}).
  \]
\end{theorem}

\begin{proof}
  \uses{thm:dr_existence, def:pixton_formula, prop:pixton_polynomiality,
        lem:invariance_weight_translation, lem:invariance_unweighted_marking}
  By the weight-translation and unweighted-marking invariances of both sides (see
  \cref{lem:invariance_weight_translation,lem:invariance_unweighted_marking}),
  it suffices to prove the case $d = 0$, $n = 0$.  In that case, the DR cycle is
  realized as the Abel-Jacobi image in an ambient $\mathbb{P}^N$ via the formula
  for DR cycles on target varieties (Janda--Pandharipande--Pixton--Zvonkine 2018);
  taking $N \to \infty$ and comparing with Pixton's formula gives the identity.
\end{proof}

\subsection{Vanishing}

\begin{theorem}[Vanishing]
  \label{thm:vanishing}
  \lean{MR4628606.TODO.vanishing}
  \uses{thm:main, def:pixton_formula}
  % SOURCE: [Bae--Holmes--Pandharipande--Schmitt--Schwarz], Theorem 3 (read from references/MR4628606-picard-stack.tex)
  % SOURCE QUOTE: "Let $g\geq 0$ and $d\in \mathbb{Z}$.
% Let $A=(a_1,\ldots,a_n)$ be a vector of integers with
% $\sum_{i=1}^n a_i=d$.
%  Pixton's vanishing holds
%  in operational Chow:
% $$ \P_{g,A,d}^c\ = 0 \
% \in {\mathsf{CH}}^c_{\mathsf{op}}(\Picabs_{g,n,d})\, \ \ \
% {\text{for all}} \ \ c>g\, ."
  \textit{Source: Bae--Holmes--Pandharipande--Schmitt--Schwarz, Theorem 3.}
  Let $g \geq 0$, $d \in \mathbb{Z}$, and $A = (a_1,\ldots,a_n)$ with $\sum a_i = d$.
  Pixton's vanishing holds in operational Chow:
  \[
    \Pix^c_{g,A,d} = 0 \in \CHop^c(\Picabs_{g,n,d}) \quad \text{for all } c > g.
  \]
\end{theorem}

\begin{proof}
  \uses{thm:main, def:pixton_formula}
  By \cref{thm:main}, $\Pix^g_{g,A,d} = \DRop_{g,A}$ is an operational class of
  codimension exactly $g$.  The classes $\Pix^c_{g,A,d}$ for $c > g$ vanish by the
  degree-$g$ property of the DR cycle's geometric origin, using the vanishings of
  Clader--Janda and Bae.
\end{proof}

\subsection{Twisted holomorphic and meromorphic differentials}

\begin{theorem}[Twisted differentials: Farkas--Pandharipande Conjecture A]
  \label{thm:twisted_differentials}
  \lean{MR4628606.TODO.twistedDifferentials}
  \uses{thm:main, thm:dr_existence, def:pixton_formula}
  % SOURCE: [Bae--Holmes--Pandharipande--Schmitt--Schwarz], Theorem 4 (read from references/MR4628606-picard-stack.tex)
  % SOURCE QUOTE: "In the strictly meromorphic
% case,
% $$\mathsf{H}_{g,A} = \mathsf{P}_{g,A,2g-2}^g[\oM_{g,n}]$$
% for the universal family
% $$\pi: \mathcal{C}_{g,n} \rightarrow \overline{\cM}_{g,n}\, ,\ \ \ \
% \omega_\pi \rightarrow \mathcal{C}_{g,n}\,. $$"
  \textit{Source: Bae--Holmes--Pandharipande--Schmitt--Schwarz, Theorem 4.}
  In the strictly meromorphic case, where $A = (a_1,\ldots,a_n)$ contains at least one
  strictly negative part and $\sum a_i = 2g-2$, the weighted fundamental class
  $\mathsf{H}_{g,A} \in \Chow_{2g-3+n}(\oM_{g,n})$ of the compact moduli space of
  twisted meromorphic differentials satisfies
  \[
    \mathsf{H}_{g,A} = \Pix^g_{g,A,2g-2}[\oM_{g,n}],
  \]
  where $\Pix^g_{g,A,2g-2}[\oM_{g,n}]$ denotes the action of Pixton's operational
  class on $[\oM_{g,n}]$ via the map $\phi_\omega\colon \oM_{g,n} \to \Picabs_{g,n,2g-2}$
  induced by the relative dualizing sheaf.
  This is Conjecture A of Farkas--Pandharipande.
\end{theorem}

\begin{proof}
  \uses{thm:main, thm:dr_existence}
  The map $\phi_\omega\colon \oM_{g,n} \to \Picabs_{g,n,2g-2}$ is induced by the universal
  family $(\pi\colon \mathcal{C}_{g,n} \to \oM_{g,n},\; \omega_\pi \to \mathcal{C}_{g,n})$.
  By Holmes's infinitesimal criterion and the compatibility of \cref{thm:dr_existence},
  $\DRop_{g,A}(\phi_\omega)([\oM_{g,n}]) = \mathsf{H}_{g,A}$.
  Applying \cref{thm:main} gives $\mathsf{H}_{g,A} = \Pix^g_{g,A,2g-2}[\oM_{g,n}]$.
\end{proof}

\subsection{Invariance properties}

\begin{lemma}[Dualizing invariance]
  \label{lem:invariance_dualizing}
  \lean{MR4628606.TODO.invarianceDualizing}
  \uses{thm:dr_existence}
  % SOURCE: [Bae--Holmes--Pandharipande--Schmitt--Schwarz], Section 0.3, Invariance I (read from references/MR4628606-picard-stack.tex)
  % SOURCE QUOTE: "Let $\mathsf{DR}^{\mathsf{op}}_{g,-A, \mathcal{L}^*}\in \mathsf{CH}^g_{\mathsf{op}}(\mathcal{S})$ be the twisted double ramification cycle associated to the new family \eqref{ff44992} and the vector $-A=(-a_1,\ldots,-a_n)$. We have the invariance
%         $$\mathsf{DR}^{\mathsf{op}}_{g,-A, \mathcal{L}^*}
%         =\epsilon^* \mathsf{DR}^{\mathsf{op}}_{g,A, \mathcal{L}}
%         \, , $$
% where $\epsilon: \Picabs_{g,n,-d} \rightarrow \Picabs_{g,n,d}$
% is the natural map obtained via dualizing the line bundle."
  \textit{Source: Bae--Holmes--Pandharipande--Schmitt--Schwarz, Section 0.3, Invariance I.}
  Let $\epsilon\colon \Picabs_{g,n,-d} \to \Picabs_{g,n,d}$ be the map obtained by dualizing
  the line bundle.  Then
  \[
    \DRop_{g,-A,\mathcal{L}^*} = \epsilon^* \DRop_{g,A,\mathcal{L}}.
  \]
\end{lemma}

\begin{proof}
  \uses{thm:dr_existence}
  The log-geometric construction of $\DRop_{g,A,\mathcal{L}}$ is symmetric under replacing
  $\mathcal{L}$ by its dual $\mathcal{L}^*$ and $A$ by $-A$: the map
  $\epsilon$ dualizes the universal line bundle, and the Abel-Jacobi compactification is
  compatible with this symmetry. Explicitly, for any test family $\phi: S \to \Picabs_{g,n,d}$
  with line bundle $\mathcal{L}$, the class $\epsilon^*\DRop_{g,A,\mathcal{L}}(\phi)$ agrees
  with $\DRop_{g,-A,\mathcal{L}^*}(\epsilon \circ \phi)$ by the functoriality of pullback
  and the definition of the DR cycle via log compactification.
\end{proof}

\begin{lemma}[Unweighted marking invariance]
  \label{lem:invariance_unweighted_marking}
  \lean{MR4628606.TODO.invarianceUnweightedMarking}
  \uses{thm:dr_existence}
  % SOURCE: [Bae--Holmes--Pandharipande--Schmitt--Schwarz], Section 0.3, Invariance II (read from references/MR4628606-picard-stack.tex)
  % SOURCE QUOTE: "Let $A_0 \in \mathbb{Z}^{n+1}$ be the vector obtained by appending $0$ (as the last coefficient) to $A$.
%  Let $\mathsf{DR}^{\mathsf{op}}_{g,A_0, \mathcal{L}}\in \mathsf{CH}^g_{\mathsf{op}}(\mathcal{S})$ be the twisted double ramification cycle associated to the new family \eqref{ff44993} and the vector $A_0$. We have the invariance
%         $$\mathsf{DR}^{\mathsf{op}}_{g,A_0, \mathcal{L}}
%         =F^*\mathsf{DR}^{\mathsf{op}}_{g,A, \mathcal{L}}
%         \, , $$
%     where $F:\Picabs_{g,n+1,d} \rightarrow \Picabs_{g,n,d}$
%     is the map forgetting the last marking."
  \textit{Source: Bae--Holmes--Pandharipande--Schmitt--Schwarz, Section 0.3, Invariance II.}
  Let $A_0 = (a_1,\ldots,a_n,0) \in \mathbb{Z}^{n+1}$ and
  $F\colon \Picabs_{g,n+1,d} \to \Picabs_{g,n,d}$ the map forgetting the last marking.
  Then
  \[
    \DRop_{g,A_0,\mathcal{L}} = F^* \DRop_{g,A,\mathcal{L}}.
  \]
\end{lemma}

\begin{proof}
  \uses{thm:dr_existence}
  Adding a marking with weight $0$ does not change the locus where
  $\cO_C(\sum a_i p_i) \simeq \mathcal{L}|_C$, since the condition on the new section $p_{n+1}$
  is vacuous.  In the log-geometric construction, the stack of tropical divisors for $A_0$
  is the pullback via $F$ of the stack for $A$, so the DR cycle pulls back accordingly.
  More precisely, for any test family $\phi: S \to \Picabs_{g,n+1,d}$, the class
  $F^*\DRop_{g,A,\mathcal{L}}(\phi)$ is computed via $F\circ\phi: S\to\Picabs_{g,n,d}$
  and coincides with $\DRop_{g,A_0,\mathcal{L}}(\phi)$ by the universal property of $F$.
\end{proof}

\begin{lemma}[Weight translation invariance]
  \label{lem:invariance_weight_translation}
  \lean{MR4628606.TODO.invarianceWeightTranslation}
  \uses{thm:dr_existence}
  % SOURCE: [Bae--Holmes--Pandharipande--Schmitt--Schwarz], Section 0.3, Invariance III (read from references/MR4628606-picard-stack.tex)
  % SOURCE QUOTE: "Let $B=(b_1,\ldots,b_n) \in \bb Z^n$ satisfy $\sum_{i=1}^n b_i=e$, then the family
%         \begin{equation}
%         ...
%         \mathcal{L}\big(\sum_{i=1}^n b_i p_i\big) \rightarrow \mathcal{C}\, .
%         \end{equation}
%  defines an object of $\Picabs_{g,n,d+e}$.
%  Let $\mathsf{DR}^{\mathsf{op}}_{g,A+B, \mathcal{L}(\sum_i b_i p_i)}\in \mathsf{CH}^g_{\mathsf{op}}(\mathcal{S})$ be the twisted double ramification cycle associated to the new family \eqref{ff44994} and the vector $A+B$. We have the invariance
%         $$\mathsf{DR}^{\mathsf{op}}_{g,A+B, \mathcal{L}(\sum_i b_i p_i)}
%         =\mathsf{DR}^{\mathsf{op}}_{g,A, \mathcal{L}}
%         \, . $$"
  \textit{Source: Bae--Holmes--Pandharipande--Schmitt--Schwarz, Section 0.3, Invariance III.}
  For $B = (b_1,\ldots,b_n) \in \mathbb{Z}^n$ with $\sum b_i = e$,
  \[
    \DRop_{g,A+B,\,\mathcal{L}(\sum_i b_i p_i)} = \DRop_{g,A,\mathcal{L}}.
  \]
\end{lemma}

\begin{proof}
  \uses{thm:dr_existence}
  Replacing $\mathcal{L}$ by $\mathcal{L}(\sum_i b_i p_i)$ and $A$ by $A+B$ does not
  change the set-theoretic locus where $\cO_C(\sum_i(a_i+b_i)p_i)\simeq\mathcal{L}(\sum_i b_i p_i)|_C$,
  since both conditions are equivalent to $\cO_C(\sum_i a_i p_i)\simeq\mathcal{L}|_C$.
  The Abel-Jacobi map and its log-geometric compactification transform compatibly under
  this twist, yielding the identity of operational classes.
\end{proof}

\begin{lemma}[Twisting by pullback]
  \label{lem:invariance_twisting_pullback}
  \lean{MR4628606.TODO.invarianceTwistingPullback}
  \uses{thm:dr_existence}
  % SOURCE: [Bae--Holmes--Pandharipande--Schmitt--Schwarz], Section 0.3, Invariance IV (read from references/MR4628606-picard-stack.tex)
  % SOURCE QUOTE: "Let $\mathcal{B} \rightarrow \mathcal{S}$ be any line bundle on the
% base. By tensoring \eqref{ff449911} with $\pi^*\mathcal{B}$, we obtain a new object of $\Picabs_{g,n,d}$ over $\mathcal{S}$:
% \begin{equation}
% ...
%         \mathcal{L}\otimes \pi^*\mathcal{B} \rightarrow \mathcal{C}\, .
%         \end{equation}
%         Let $\mathsf{DR}^{\mathsf{op}}_{g,A, \mathcal{L}\otimes\pi^*\mathcal{B}}\in \mathsf{CH}^g_{\mathsf{op}}(\mathcal{S})$ be the twisted double ramification cycle associated to the new family \eqref{ff4499} and the vector $A$. We have the invariance
%         $$\mathsf{DR}^{\mathsf{op}}_{g,A, \mathcal{L}\otimes\pi^*\mathcal{B}}
%         =\mathsf{DR}^{\mathsf{op}}_{g,A, \mathcal{L}}
%         \, . $$"
  \textit{Source: Bae--Holmes--Pandharipande--Schmitt--Schwarz, Section 0.3, Invariance IV.}
  Let $\mathcal{B} \to \mathcal{S}$ be any line bundle on the base, and form the new object
  of $\Picabs_{g,n,d}$ obtained by tensoring the universal line bundle with $\pi^*\mathcal{B}$,
  i.e.\ the family with line bundle $\mathcal{L} \otimes \pi^*\mathcal{B} \to \mathcal{C}$.
  Then
  \[
    \DRop_{g,A,\,\mathcal{L} \otimes \pi^*\mathcal{B}} = \DRop_{g,A,\mathcal{L}}.
  \]
\end{lemma}

\begin{proof}
  \uses{thm:dr_existence}
  Tensoring the line bundle $\mathcal{L}$ by the pullback $\pi^*\mathcal{B}$ of a line bundle
  from the base does not change the relative Abel-Jacobi data along the fibres of $\pi$:
  the comparison $\cO_C(\sum_i a_i p_i) \simeq \mathcal{L}|_C$ on each fibre is unaffected by
  twisting with a bundle pulled back from $\mathcal{S}$.  Hence the log-geometric
  compactification, and the resulting operational class, are unchanged.
\end{proof}

\begin{lemma}[Vertical twisting]
  \label{lem:invariance_vertical_twisting}
  \lean{MR4628606.TODO.invarianceVerticalTwisting}
  \uses{thm:dr_existence, lem:invariance_twisting_pullback}
  % SOURCE: [Bae--Holmes--Pandharipande--Schmitt--Schwarz], Section 0.3, Invariance V (read from references/MR4628606-picard-stack.tex)
  % SOURCE QUOTE: "Consider a partition of the genus, marking, and degree data,
% \begin{equation}\label{ss445}
% g_1+g_2=g\, , \ \ \ N_1\sqcup N_2 = \{1,\ldots,n\}\,, \ \ \ d_1+d_2=d\, ,
% \end{equation}
% which is {\em not} symmetric. ...
% Such a partition defines a divisor
% $$ \Delta_1 \in \mathsf{CH}^1(\mathfrak{C}_{g,n,d})$$
% in the universal curve over $\Picabs_{g,n,d}$ ...
% By tensoring \eqref{ff449911} with $\mathcal{O}_{\mathcal{C}}(\Delta_1)$, we obtain a new object of $\Picabs_{g,n,d}$ over $\mathcal{S}$: ...
%         \mathcal{L}(\Delta_1) \rightarrow \mathcal{C}\, . ...
%         Let $\mathsf{DR}^{\mathsf{op}}_{g,A, \mathcal{L}(\Delta_1)}\in \mathsf{CH}^g_{\mathsf{op}}(\mathcal{S})$ be the twisted double ramification cycle associated to the new family \eqref{ff44999} and the vector $A$. We have the invariance
%         \begin{equation}\label{k55k9}
%         \mathsf{DR}^{\mathsf{op}}_{g,A, \mathcal{L}(\Delta_1)} =
%         \mathsf{DR}^{\mathsf{op}}_{g,A, \mathcal{L}}
%         \, . \end{equation}"
  \textit{Source: Bae--Holmes--Pandharipande--Schmitt--Schwarz, Section 0.3, Invariance V.}
  Consider a partition of the genus, marking, and degree data
  \[
    g_1 + g_2 = g, \qquad N_1 \sqcup N_2 = \{1,\ldots,n\}, \qquad d_1 + d_2 = d,
  \]
  which is \emph{not} symmetric (i.e.\ $(g_1,N_1,d_1) \neq (g_2,N_2,d_2)$).
  Such a partition defines a divisor $\Delta_1 \in \Chow^1(\mathfrak{C}_{g,n,d})$ in the
  universal curve over $\Picabs_{g,n,d}$, by twisting by the $(g_1,N_1,d_1)$-component of a
  curve with a separating node carrying the above separating data.
  Forming the new object of $\Picabs_{g,n,d}$ with line bundle
  $\mathcal{L}(\Delta_1) \to \mathcal{C}$, we have the invariance
  \[
    \DRop_{g,A,\,\mathcal{L}(\Delta_1)} = \DRop_{g,A,\mathcal{L}}.
  \]
\end{lemma}

\begin{proof}
  \uses{thm:dr_existence, lem:invariance_twisting_pullback}
  For symmetric separating data, taking $\Delta_1 \subset \mathfrak{C}_{g,n,d}$ to be the
  full preimage of the locus $\Delta \subset \Picabs_{g,n,d}$ of curves with a separating
  node, the equality follows from \cref{lem:invariance_twisting_pullback} with
  $\mathcal{B} = \cO(\Delta)$.  In the general (non-symmetric) case, the twist by the divisor
  $\Delta_1$ likewise does not change the relative Abel-Jacobi locus, and the log-geometric
  compactification produces the same operational class.
\end{proof}

\begin{lemma}[Partial stabilization]
  \label{lem:invariance_partial_stabilization}
  \lean{MR4628606.TODO.invariancePartialStabilization}
  \uses{thm:dr_existence}
  % SOURCE: [Bae--Holmes--Pandharipande--Schmitt--Schwarz], Section 0.3, Invariance VI (read from references/MR4628606-picard-stack.tex)
  % SOURCE QUOTE: "Consider a second family
% of prestable
% $n$-pointed genus $g$ curves over $\mathcal{S}$,
% ...
% together with a birational $\mathcal{S}$-morphism
% $$f: \mathcal{C}' \rightarrow \mathcal{C}\, , \ \ \ \ f\circ p'_i = p_i\, .$$
% A line bundle of relative degree $d$ is defined on $\mathcal{C}'$ by
% $$f^*\mathcal{L} \rightarrow \mathcal{C}'\, .$$
% We require the following property to hold: {\em if the section $p'_i$ meets the
% exceptional locus of $f$, then $a_i=0$}.
%         Let $\mathsf{DR}^{\mathsf{op}}_{g,A, f^*\mathcal{L}}\in \mathsf{CH}^g_{\mathsf{op}}(\mathcal{S})$ be the twisted double ramification cycle associated to the new family ... and the vector $A$.
%         We have the invariance
%         $$\mathsf{DR}^{\mathsf{op}}_{g,A, f^*\mathcal{L}}=
%         \mathsf{DR}^{\mathsf{op}}_{g,A, \mathcal{L}}\, . $$"
  \textit{Source: Bae--Holmes--Pandharipande--Schmitt--Schwarz, Section 0.3, Invariance VI.}
  Consider a second family $\pi'\colon \mathcal{C}' \to \mathcal{S}$ of prestable
  $n$-pointed genus-$g$ curves with sections $p'_1,\ldots,p'_n$, together with a birational
  $\mathcal{S}$-morphism $f\colon \mathcal{C}' \to \mathcal{C}$ satisfying $f \circ p'_i = p_i$.
  Equip $\mathcal{C}'$ with the line bundle $f^*\mathcal{L}$ of relative degree $d$, and
  require that \emph{if the section $p'_i$ meets the exceptional locus of $f$, then $a_i = 0$}.
  Then
  \[
    \DRop_{g,A,\,f^*\mathcal{L}} = \DRop_{g,A,\mathcal{L}}.
  \]
\end{lemma}

\begin{proof}
  \uses{thm:dr_existence}
  The birational morphism $f$ contracts the exceptional locus, and the pulled-back bundle
  $f^*\mathcal{L}$ has the same relative degree $d$.  Since the marked points meeting the
  exceptional locus carry weight $a_i = 0$, the Abel-Jacobi comparison
  $\cO_C(\sum_i a_i p_i) \simeq \mathcal{L}|_C$ is preserved under pullback along $f$.
  Thus the log-geometric construction yields the same universal twisted double ramification
  cycle on the base $\mathcal{S}$.
\end{proof}

\subsection{Universal formula in degree 0}

\begin{theorem}[Universal formula in degree 0]
  \label{thm:universal_deg0}
  \lean{MR4628606.TODO.universalDeg0}
  \uses{thm:main, def:pixton_formula, def:tautological_classes, def:operational_chow,
        def:picard_stack, def:boundary_map_j, lem:invariance_unweighted_marking,
        lem:invariance_dualizing}
  % SOURCE: [Bae--Holmes--Pandharipande--Schmitt--Schwarz], Section 0.7 (sec:intro_deg_0) (read from references/MR4628606-picard-stack.tex)
  % SOURCE QUOTE: "The specialization of Theorem \ref{Thm:main} to $d=0$ calculates $\DRop_{g,\emptyset}$ as
% the value at $r=0$ of the degree $g$ part of
% \begin{align*}
% \exp\left(-\frac12 \eta \right) \sum_{
% \substack{\Gamma_\delta\in \G_{g,0,0} \\
% w\in \mathsf{W}_{\Gamma_\delta,r}}
% }
% \frac{r^{-h^1(\Gamma_\delta)}}{|\Aut(\Gamma_\delta)| }
% \;
% j_{\Gamma_\delta*}\Bigg[&
% \prod_{e=(h,h')\in \E(\Gamma_\delta)}
% \frac{1-\exp\left(-\frac{w(h)w(h')}2(\psi_h+\psi_{h'})\right)}{\psi_h + \psi_{h'}} \Bigg]\,
% \end{align*}
% as an operational Chow class on
% $$\Picabs_g=\Picabs_{g,0,0}\, .$$ ...
% The full
% statement of Theorem \ref{Thm:main} can
% be recovered from the above $d=0$ specialization via pullback under the map
% \[\Picabs_{g,n,d} \to \Picabs_{g}\, ,\ \ \  (C,p_1, \ldots, p_n, \ca L) \mapsto \left(C, \ca L\left(-\sum_{i=1}^n a_i p_i\right)\right).\]"
  % SOURCE QUOTE (compact type): "By restricting to suitable open subsets of $\Picabs_{g}$, we can simplify the $d=0$ formula even further. Let $$\Picabs_{g}^\textup{ct} \subset \Picabs_{g}$$
% be the locus where the curve $C$ is of compact type.
% We obtain
% \begin{equation}\label{ccttt}
%  \DRop_{g,\emptyset}|_{\Picabs_{g}^\textup{ct}}
%  = \frac{\theta^g}{g!} \,, \ \ \
%  \text{for}\ \ \theta = -\frac12 \left(\eta +  \sum_{\Delta} d_\Delta^2 [\Delta] \right),
% \end{equation}
% where the sum is over the boundary divisors $\Delta \subset \Picabs_g$ on which generically the curve splits into two components carrying line bundles of degrees $d_\Delta$ and $-d_\Delta$."
  \textit{Source: Bae--Holmes--Pandharipande--Schmitt--Schwarz, Section 0.7.}
  Let $g \neq 1$.  The specialization of \cref{thm:main} to $d = 0$ with no markings
  computes $\DRop_{g,\emptyset} \in \CHop^g(\Picabs_g)$, where $\Picabs_g = \Picabs_{g,0,0}$,
  as the value at $r = 0$ of the degree-$g$ part of
  \[
    e^{-\frac{1}{2}\eta}
    \sum_{\substack{\Gamma_\delta \in \GG_{g,0,0} \\ w \in \mathsf{W}_{\Gamma_\delta,r}}}
    \frac{r^{-h^1(\Gamma_\delta)}}{|\Aut(\Gamma_\delta)|}
    \; j_{\Gamma_\delta*}\!\left[
      \prod_{e=(h,h') \in \E(\Gamma_\delta)}
        \frac{1 - e^{-\frac{w(h)w(h')}{2}(\psi_h + \psi_{h'})}}{\psi_h + \psi_{h'}}
    \right].
  \]
  The full statement of \cref{thm:main} for general $A$ and $d$ is recovered from this
  $d = 0$, $A = \emptyset$ specialization via pullback under the map
  \[
    \Picabs_{g,n,d} \to \Picabs_g, \qquad
    (C, p_1,\ldots,p_n, \mathcal{L}) \mapsto \Big(C, \mathcal{L}\big(-\textstyle\sum_{i=1}^n a_i p_i\big)\Big),
  \]
  which factors as $\tau_{-A}\colon \Picabs_{g,n,d} \to \Picabs_{g,n,0}$ followed by the
  marking-forgetting map $F\colon \Picabs_{g,n,0} \to \Picabs_g$.
  Restricting to the compact-type locus
  $\Picabs_g^{\mathrm{ct}} \subset \Picabs_g$ (where $C$ is of compact type), the formula
  simplifies to the theta-divisor expression
  \[
    \DRop_{g,\emptyset}\big|_{\Picabs_g^{\mathrm{ct}}} = \frac{\theta^g}{g!},
    \qquad
    \theta = -\frac{1}{2}\Big(\eta + \sum_{\Delta} d_\Delta^2\, [\Delta]\Big),
  \]
  where the sum is over the boundary divisors $\Delta \subset \Picabs_g$ on which the curve
  generically splits into two components carrying line bundles of degrees $d_\Delta$ and
  $-d_\Delta$.  The class $\theta$ is the universal theta divisor on
  $\Picabs_g^{\mathrm{ct}}$.
\end{theorem}

\begin{proof}
  \uses{thm:main, def:pixton_formula, lem:invariance_unweighted_marking, lem:invariance_dualizing}
  The displayed degree-$0$ formula is the specialization of the definition of
  $\Pix^g_{g,\emptyset,0}$ to $\GG_{g,0,0}$, where the leg factors are absent and the vertex
  factors $\prod_v e^{-\frac{1}{2}\eta(v)}$ collapse to the global factor $e^{-\frac{1}{2}\eta}$;
  the equality with $\DRop_{g,\emptyset}$ is \cref{thm:main}.  For the reconstruction of the
  general case, the unweighted-marking invariance (\cref{lem:invariance_unweighted_marking})
  gives $F^*\DRop_{g,\emptyset} = \DRop_{g,\mathbf{0}}$ for the zero vector
  $\mathbf{0} \in \mathbb{Z}^n$, and the dualizing/weight-translation invariance applied to the
  map $\tau_{-A}$ (\cref{lem:invariance_dualizing}) gives
  $\tau_{-A}^*\DRop_{g,\mathbf{0}} = \DRop_{g,A}$; composing yields the general formula, and
  the parallel symmetries of Pixton's formula give $\tau_{-A}^*F^*\Pix^g_{g,\emptyset,0} =
  \Pix^g_{g,A,d}$.  Over the compact-type locus, every $h^1$ contribution comes from
  separating nodes, so the edge factors reduce to powers of the divisor classes $[\Delta]$
  and the formula collapses to $\theta^g/g!$ with $\theta$ as stated.
\end{proof}
